%% kapa: exam=1 round=21 year=2010 period=1 form=A subject=civil_law subject_ko=민법 seq=38 number=38
%% source: Q-Net 감정평가사 기출문제 (cst00309) artlSeq=5024181 file=extracted/1008/5024181/2010년 제21회 감정평가사 1차시험 기출문제/1교시.hwp
%% conversion: kordoc 4.24.51 -> tools/gen1.py (auto); review=proofread-sonnet; flags=none
\begin{question}{38}
가등기담보 등에 관한 법률에 관한 설명으로 옳은 것을 모두 고른 것은? (다툼이 있으면 판례에 의함)
\begin{framedtext}
ㄱ. 가등기담보 등에 관한 법률은 담보가등기권자에게 경매청구권과 경매절차에서의 우선변제권을 부여하는 등 저당권자와 유사한 지위를 인정한다. \par
ㄴ. 귀속청산의 경우 변제기가 경과하더라도 담보가등기권자가 곧바로 담보권을 실행할 수는 없고 변제기 후 일정한 기간이 경과할 것을 요한다. \par
ㄷ. 귀속청산의 경우 후순위권리자는 담보가등기권자가 채무자에게 지급하여야 할 청산금에 대하여 아무런 권리도 행사할 수 없다. \par
ㄹ. 가등기담보의 성질상 가등기담보권의 양도는 허용되지 않는다. \par
ㅁ. 귀속청산의 경우 담보가등기권자가 채무자에게 지급할 청산금이 없는 때에는 2월의 청산기간이 경과함으로써 청산절차는 종료되고, 이후 담보목적물에 대한 사용ㆍ수익권은 담보가등기권자에게 귀속된다. \par
\end{framedtext}
\begin{choices}
  \item ㄱ, ㄹ
  \item ㄴ, ㄷ
  \item ㄱ, ㄴ, ㅁ
  \item ㄷ, ㄹ, ㅁ
  \item ㄱ, ㄴ, ㄹ, ㅁ
\end{choices}
\end{question}
\answer{3}
